\scanpage{149}\mspage{4}%
% The first occurrence of strictes has one wavy underline; the second has one nearly straight underline.
pé en long et en large dans le livre de Mme Sinh, je pense, sous le titre de chapitre des ``Gr-catégories \uline{strictes}'' (i.e. \inserted{\unsure{où}} les isomorphismes d'associativité, d'unité, et d'inverse $X\mathbin{.}X^{-1}\unsure{\simeq|=}1$ sont des \uline{identités}). On peut expliciter p.ex. comment (via le fait que tte Gr-cat. est Gr-équivalente à une Gr-catégorie stricte) le calcul avec les Gr-catégories se \unsure{réduit} à un calcul très algébrique avec des Gr-catégories \uline{strictes}, par une sorte de ``calcul des fractions'' (à droite ou à gauche, au choix) du type de celui qui aboutit à la construction des catégories dérivées. Quoi qu'il en soit, voici l'explicitation de cette structure (groupe dans $(1\text{-ct.})$) en termes de théorie des groupes ($1$-catégorie dans $(\mathrm{Gr})$) :

% The isolated low mark immediately before quadruplet may be a pen-contact mark rather than punctuation.
La structure est décrite par un \unsure{,|.}quadruplet $(L_1,L_0,d,\unsure{\Theta|\theta})$, avec
% The two small marks to the right of L_0 may be a colon or surface/pen specks.
\diagram{\begin{tikzcd}[ampersand replacement=\&,column sep=3em]
L_1 \arrow[r,"d"] \& L_0\;\unsure{:}
\end{tikzcd}}
un hom de groupes ordinaires, et
\[
\Theta:L_0\longrightarrow\underline{\mathrm{Aut}}_{\mathrm{gr}}(L_1)
\]
une opération de $L_0\deleted{\unsure{\Theta|0}}$ sur $L_1$, avec les
% archive sheet number: 149
